% ============================================================
% 级联错误传播假设 —— 修订稿
% 用于替换 chap03.tex 第208-269行的内容
% ============================================================

\subsection{级联错误传播假设}
\label{subsec:cascading_error}

在多智能体协作系统中，错误信息的传播机制是理解系统安全性的关键问题。本文提出级联错误传播假设，其核心观点为：在链式连接的多智能体系统中，早期智能体产生的微小错误会在下游传播过程中被逐步放大，导致最终输出与正确方案的偏离程度显著大于初始错误的幅度。

该假设的提出源于对大语言模型上下文依赖行为的观察。现代大语言模型采用自回归生成方式，每个输出token的预测都依赖于之前的所有内容。在多智能体协作场景中，下游智能体的输入包含了上游智能体的输出结果，上游输出中的任何错误信息都会成为下游推理的前提。大语言模型倾向于保持与输入上下文的语义一致性，这种特性在正常场景下有助于生成连贯的对话，但在错误传播场景下却产生负面效应。当上游智能体输出了包含错误医学知识的内容后，下游智能体在接收该上下文时，会倾向于顺从上游的判断而非主动纠正错误。本文将这种现象定义为"上下文锚定效应"（Context Anchoring Effect），其机理类似于认知科学中的锚定偏见\cite{tversky1974judgment}：下游智能体将上游输出视为已经过验证的可靠信息，在此基础上进行后续推理，即使该智能体本身具备正确的医学知识，也难以完全克服错误上下文的影响。

上述定性分析可以进一步形式化。根据第\ref{subsec:state_transition}节的状态转移方程，智能体$A_i$的上下文仅由原始查询$Q$和直接前驱$A_{i-1}$的输出$O_{i-1}$构成，即$Context_i = \langle Q, O_{i-1} \rangle$。这意味着$A_i$的输出状态仅直接取决于$A_{i-1}$的输出状态，与更早的智能体$A_1, \ldots, A_{i-2}$之间不存在直接的信息传递通道。换言之，错误信息在链中的传播满足马尔可夫性：给定$A_{i-1}$的输出状态，$A_i$的错误概率与$A_1, \ldots, A_{i-2}$的状态条件独立。

基于这一马尔可夫性质，定义事件$E_i$表示第$i$个智能体输出包含错误医疗信息，其对立事件$\bar{E}_i$表示正确输出。根据全概率公式，下游智能体$A_i$的错误概率可按上游智能体$A_{i-1}$的错误状态进行条件分解：

\begin{equation}
P(E_i) = P(E_i | E_{i-1}) \cdot P(E_{i-1}) + P(E_i | \bar{E}_{i-1}) \cdot P(\bar{E}_{i-1})
\label{eq:error_propagation}
\end{equation}

该公式表明，下游智能体的错误率由两部分组成：上游已发生错误时的错误传播概率，以及上游正常时自身产生错误的概率。由于上下文锚定效应的存在，通常有$P(E_i | E_{i-1}) > P(E_i | \bar{E}_{i-1})$，即上游错误的存在会显著提高下游犯错的概率。

在此基础上，为了量化上游错误对下游智能体的影响程度，本文定义错误放大因子（Error Amplification Factor, $\alpha$）。设智能体链为$A = \{A_1, A_2, \ldots, A_n\}$，第$i$个位置的错误放大因子定义为：

\begin{equation}
\alpha_i = \frac{P(E_{i+1} | E_i)}{P(E_{i+1} | \bar{E}_i)}
\label{eq:error_amplification}
\end{equation}

$\alpha_i$度量的是上游错误相对于上游正确时，下游智能体犯错概率的倍增程度。当$\alpha_i > 1$时，上游错误增加了下游犯错的可能性，错误被放大传播；当$\alpha_i = 1$时，上下游错误状态相互独立，不存在传播效应；当$\alpha_i < 1$时，下游智能体在一定程度上起到了纠错作用。从形式上看，$\alpha_i$与流行病学中衡量暴露因素与结果事件关联强度的相对风险（Relative Risk）指标\cite{rothman2008epidemiology}具有相同的数学结构，本研究将这一成熟的统计概念引入多智能体系统分析，用于量化级联效应的强度。

由于错误传播满足马尔可夫性，可以推导出链式结构中首个智能体的错误状态对任意下游智能体的影响。考虑从$A_1$到$A_n$的完整链路，$A_1$发生错误时$A_n$的条件错误概率与$A_1$正常时$A_n$的条件错误概率之比，可定义为总放大因子$\alpha_{total}$。当各相邻位置之间的错误传播相互独立（即$\alpha_i$不随上游错误的具体传播路径而改变）时，总放大因子等于各位置放大因子的乘积：

\begin{equation}
\alpha_{total} = \prod_{i=1}^{n-1} \alpha_i
\label{eq:total_amplification}
\end{equation}

该公式揭示了链式结构长度对系统安全性的重要影响。若每个位置的$\alpha_i$均大于1，则随着链长度$n$的增加，$\alpha_{total}$将呈指数增长，表明上游错误对下游的影响迅速放大。需要指出的是，$\alpha_{total}$本身并非概率值，而是条件概率的比值，因此不能由此直接推断"最终输出错误的概率趋近于1"。但高$\alpha_{total}$意味着，一旦首个智能体输出错误，链尾智能体犯错的概率将远高于首个智能体正确时的水平。从防御角度看，这一结论指出了保护首个智能体的关键意义：若$A_1$能够抵御攻击并保持正确输出，即使后续链路存在一定程度的脆弱性，整体错误风险仍可控制在较低水平。

基于上述分析，本文提出多智能体系统的安全评估应重点关注两个维度。一是位置脆弱性，即评估不同位置的智能体被攻击后对系统整体的影响差异，由于链首位置的错误能够波及所有下游节点，该位置通常具有最高的脆弱性。二是放大强度，即通过在实验中计算实际的错误放大因子来验证上述理论预测的正确性。这两个维度的分析为多智能体医疗协作系统的安全评估提供了理论框架，也为第\ref{cha:multi_agent}章的实验设计明确了核心假设。

此外需要说明的是，上述分析依赖的"盲目信任"前提，即下游智能体将上游输出作为可信上下文、不会主动质疑或验证其正确性，是基于当前多智能体框架的默认行为。AutoGen、LangChain等主流框架均采用类似设计，以避免过度验证带来的响应延迟。揭示这一设计缺陷正是本研究的安全评估目标之一。未来具备交叉验证能力的多智能体系统有望缓解级联错误传播问题，但评估此类增强系统的有效性本身就需要以理解当前系统的脆弱性为前提。图\ref{fig:cascade_propagation}展示了错误在智能体链中的传播与放大过程。

\begin{figure}[htbp]
\centering
\fbox{\parbox{0.85\textwidth}{\centering\vspace{3cm}
\textbf{图片占位符}\\
级联错误传播与放大机制示意图\\四个智能体节点 $A_1 \to A_2 \to A_3 \to A_4$\\
错误从$A_1$传播到$A_4$的路径，每步标注$\alpha$值\\
展示累积放大效应：$\alpha_{total} = \alpha_1 \times \alpha_2 \times \alpha_3$
\vspace{3cm}}}
\caption{级联错误传播与放大机制示意图}
\label{fig:cascade_propagation}
\end{figure}
